\section{Gott entdecken}
\begin{block}[Gott entdecken] 
        Vor zwei Wochen habe ich mit der Serie \enquote{Wer ist Gott?} begonnen. Letzten Sonntag habe ich einen Vortrag über Qumran und die Funde, die dort gemacht wurden, gehalten. Auch da kam die Suveränität Gottes zum Vorschein.

        Heute will ich darüber sprechen, wie wir Gott entdecken können oder wie sich Gott uns bemerkbar macht. Diese Serie ist nicht dafür da, dass ich euch Beweise für einen Gott präsentiere. Ich kann Gott nicht beweisen. Niemand kann das. Darüber bin ich auch froh. Wenn man Gott beweisen könnte, wäre er nicht grösser als die Beweise selber. Aber Gott ist Gott. Er sagt von sich:
        \begin{bibelbox}{SCHL}{IIMos}{3:14}
            Gott sprach zu Mose: »Ich bin, der ich bin!«
        \end{bibelbox}
        \betonung[b,p]{YHWH} 
        \\

        Ich möchte mit diesem Vortrag, mehr über Gott erfahren, ihn besser kennen lernen und dadurch auch mehr Vertrauen zu ihm aufbauen. 
        
        Es ist wie mit dem Ehepartner. Einen Ehepartner lernt man nicht kennen, wenn man frisch in ihn verliebt ist. Zu dem Zeitpunkt sieht man alles rosa und sieht über vieles hinweg, was einen später stören wird. Einen Ehepartner lernt man mit der Zeit kennen, in der man mit ihm zusammen lebt. In verschiedenen Situationen, Gefühlsstimmungen, in guten und schlechten Zeiten. So ist es auch mit Gott. Zu Beginn, wenn man sich zu Gott bekehrt, da ist man euphorisch. Man will jetzt die ganze Welt evangelisieren und verändern. Je nach dem, in welchem Umfeld man sich bekehrt hat wurde gesagt, dass man Gesund wird, das Leben leichter geht, alles einfacher wird. Anderen wird gesagt, jetzt darfst nicht mehr rauchen, du musst immer beten und unbedingt jeden Sonntag in den Gottesdienst gehen, Frauen nur noch Röcke tragen.        

        Mit der Zeit fragt man sich, wo ist denn dieser Gott, für den ich das alles mache? Niemand will das Evangelium hören, meine Krankheit ist immer noch da, meine Zornausbrüche habe ich immer noch nicht im Griff, keine Lust zum Beten, oder in den Gottesdienst zu gehen. Oder ich bete doch und Gott antwortet nicht. Man fängt an, an Gott zu zweifeln. Vielleicht nicht direkt zu zweifeln, aber doch Fragen zu stellen: \enquote{Wo und was ist denn dieser Gott, zu dem ich mich bekehrt habe?} Mit diesen Zweifeln kommen zusätzlich Fragen auf: Bin ich wirklich errettet? Oder mache ich genug für Gott, für den Glauben, für die Religion?
        In Psalm 55 lesen wir das David mit den gleichen Problemen gekämpft hat:
        \begin{bibelbox}{SCHL}{Ps}{55:2.23}
        2 Schenke meinem Gebet Gehör, o Gott, und verbirg dich nicht vor meinem Flehen!
        \\
        \textbf{Dann am Ende:}
        \\
        24 Ja, du, o Gott, wirst sie in die Grube des Verderbens hinunterstoßen; die Blutgierigen und Falschen werden es nicht bis zur Hälfte ihrer Tage bringen. Ich aber vertraue auf dich!
        \end{bibelbox}        
        Jesus verspricht uns, (einer meiner lieblings Verse um meinen Weg mit Jesus zu prüfen:)
        \begin{bibelbox}{SCHL}{Matt}{11:30}
            Denn mein Joch ist sanft und meine Last ist leicht.
        \end{bibelbox}
        \begin{block}       
        Eines dürfen wir nicht vergessen. Gott ist nicht nur exklusiv für uns Christen da. 
        \begin{bibelbox}{SCHL}{ITim}{2:4}
        welcher will, dass alle Menschen gerettet werden und zur Erkenntnis der Wahrheit kommen.
        \end{bibelbox}
        \begin{bibelbox}{SCHL}{Hes}{33:11}
        Sprich zu ihnen: So wahr ich lebe, spricht GOTT, der Herr: Ich habe kein Gefallen am Tod des Gottlosen, sondern daran, dass der Gottlose umkehre von seinem Weg und lebe! Kehrt um, kehrt um von euren bösen Wegen! Warum wollt ihr sterben, o Haus Israel?
        \end{bibelbox}
        Gott ist für alle Menschen da. Gott liebt jeden Menschen genau gleich, egal ob Christ oder nicht. Darum jubelt der Himmel, wenn sich jemand zu Gott bekehrt und ihn von Herzen annimmt.
        \begin{bibelbox}{SCHL}{Lk}{15:10}
        Ich sage euch, so ist auch Freude vor den Engeln Gottes über einen Sünder, der Buße tut.
        \end{bibelbox}
        Gott freut sich mit uns die Ewigkeit mit ihm verbringen zu können. Oder liebt Gott Thomas vielleicht mehr, nur weil er Prediger ist und die Bibel auswendig kennt? 
        \end{block}
        Woran erkennen wir jetzt Gott vier Punkte. 
        \\
        \begin{block}[Die Sehnsucht]
        Dass Gott immer schon existiert hat, habe ich euch vor zwei Wochen schon gesagt, lest \bibleverse{IMos}(1:1). Nur weil Gott existiert leben wir. Nicht umgekehrt. Nicht wir haben Gott erfunden, sondern Gott hat uns nach seinem Ebenbild erschaffen. Darum haben alle Menschen eine Sehnsucht nach etwas Höherem und eine Hoffnung nach einem Leben nach dem Tod.
        \begin{bibelbox}{SCHL}{Pred}{3:11}
        Er hat alles vortrefflich gemacht zu seiner Zeit, auch die Ewigkeit hat er ihnen ins Herz gelegt — nur dass der Mensch das Werk, das Gott getan hat, nicht von Anfang bis zu Ende ergründen kann.
        \end{bibelbox}
        Tiere haben das nicht. Es gibt keine Tierart, die ein anderes Tier begräbt oder einen Gottesdienst veranstaltet. Dieses Bedürfnis hat nur der Mensch.
        \\
        \betonung[b,p]{Es gibt also eine Sehnsucht nach einem Gott, und nach einem ewigen Leben.}
        \end{block}
        \begin{block}[Die Schöpfung]
        Gott hat für uns Menschen die Erde geschaffen. Am sechsten Tag hat er die Landtiere geschaffen und nach den Landtieren den Menschen und diesem Menschen sagt Gott, dass sie sich die Erde untertan machen sollen. 
        \begin{bibelbox}{SCHL}{IMos}{1:28}
        Und Gott segnete sie; und Gott sprach zu ihnen: Seid fruchtbar und mehrt euch und füllt die Erde und macht sie euch untertan; und herrscht über die Fische im Meer und über die Vögel des Himmels und über alles Lebendige, das sich regt auf der Erde!
        \end{bibelbox}
        Es steht nicht die Erde zerstören, sondern bebauen und pflegen. Das ganze Universum, die Erde, alles da Draussen, hat Gott nur für \betonung[b,p]{DICH UND MICH} geschaffen. \bibleverse{Ps}(104:) Und nein, dafür brauchte er keine Millionen von Jahren sondern 6 Tage. Weil er wusste, dass wir Menschen plötzlich mit solchen Märchen kommen, lässt er Mose explizit aufschreiben: \enquote{Und es wurde Abend und es wurde Morgen der xte Tag.} Das steht 6mal in 1.Mos Kapitel 1 in den Versen 5.8.13.19.23.31.

        \betonung[b,p]{Als zweites also die Schöpfung}
        \end{block}
        \begin{block}[Das Wort]
        \begin{bibelbox}{SCHL}{Joh}{1:1}
        Am Anfang war das Wort, und das Wort war bei Gott und das Wort war Gott.
        \end{bibelbox}
        Gott hat uns sein Wort gegeben und Gott ist das Wort. 
        \begin{bibelbox}{SCHL}{IITim}{3:16}
        Alle Schrift ist von Gott eingegeben und nützlich zur Belehrung, zur Überführung, zur Zurechtweisung, zur Erziehung in der Gerechtigkeit
        \end{bibelbox}
        Wer letzten Sonntag hier war, weiss jetzt, welche Funde in Qumran am Toten Meer gemacht wurden. Viele dieser Funde sind mehr als 2000 Jahre alt. Die vollständige Jesajarolle \zB{} wurde um 175 \vChr{} geschrieben und dort steht das gleiche drin, wie hier in dieser deutschen modernen Bibel. Wir dürfen nicht vergessen, dass dieses Buch, das am meisten gehasste, am meisten verfolgte, am meisten untersuchte Buch der Weltgeschichte ist. Auch in der Christenverfolgung 300 \nChr{} wo die Bibel aus losen Rollen bestand, konnte sie nicht vernichtet werden. Nie konnte dieses Buch vernichtet werden. 10tausende Wissenschaftler konnten keine Fehler darin nachweisen. Auch in Nordkorea, wo die Todesstrafe oder KZ (nicht nur für den Besitzer, sondern für die ganze Familie) auf den Besitz einer Bibel steht, gibt es Bibeln. Bruder Andrew von Open Doors hat beim Bibelschmuggeln, viele Wunder erlebt.    
        \begin{itemize}
            \item Nur aus der Bibel wissen wir, wer der Schöpfer ist. 1. Mose
            \item Aus der Bibel wissen wir, wieso die Welt so ist, wie sie ist! 1. Mos 3
            \item Aus der Bibel wissen wir, was für ein Plan Gott mit uns Menschen hat! Ganze AT und NT
            \item Aus der Bibel wissen wir, dass und wie er uns retten will! Evangelien.
            \item Nur durch die Bibel wissen wir, dass er uns liebt. Ganze Schrift.
            \item Durch die Bibel wissen wir auch, dass Gott gerecht ist. Das er Gericht halten wird am Ende der Zeit.
            \item Dass er wiederkommen wird Apg 1.11
            \item Dass er eine neue Erde und einen neuen Himmel Off 21 erschaffen wird.
        \end{itemize}       
        \betonung[b,p]{Also drittens das Wort Gottes}
        \end{block}
        \begin{block}[Das Gewissen]
        Jeder Mensch kann mehr oder weniger zwischen Gut und Böse unterscheiden. 
        \begin{bibelbox}{SCHL}{Rom}{2:14}
        Wenn nämlich Heiden, die das Gesetz nicht haben, doch von Natur aus tun, was das Gesetz verlangt, so sind sie, die das Gesetz nicht haben, sich selbst ein Gesetz,
        \end{bibelbox}
        Kein Tier kann das, oder hat ein Gewissen. Oft werden Tieren menschliche Emotionen angedichtet. Die Natur zeigt uns aber, dass es anders ist. Die Tierwelt ist absolut brutal unter sich. Ein Tier tötet nicht wie der Mensch aus Gier, sondern einfach nur zum Überleben. Wir Menschen sind da anders. Wir verlieben uns, wir brauchen Gemeinschaft, wir werden zornig, wir werden enttäuscht, wir sind stolz, wir haben Angst. Wir streben nach Macht und Besitz. Wir wir beten an, darum die ganzen Religionen. Wir können Gesetze erfinden, weil wir meinen zwischen Gut und Böse unterscheiden zu können. Das Problem ist nur, dass diese Unterscheidung nur richtig funktioniert, wenn wir uns an die Richtlinen halten, die uns Gott durch die Bibel gegeben hat.
        \begin{bibelbox}{SCHL}{Ps}{119:160}
        Die Summe deines Wortes ist Wharheit, und jede Bestimmung deiner Gerechtigkeit bleibt ewiglich.
        \end{bibelbox}
        Aber die Regierungen wissen es anscheinend besser und werden uns ins Elend führen. Egal ob SVP, SP oder Grüne.

        \betonung[b,p]{Viertens also das Erkennen von Gut und Böse}
       \end{block}

       Nächstes mal wollen wir uns mit den Eigenschaften von Gott beschäftigen. Bis dahin, sucht Gott in der Schöpfung, im Wort und im Herzen. Er wird sich euch zeigen. Danke fürs zuhören.
        \begin{enumerate}
            \item Warum der Gott der Bibel und kein anderer?
            \item Gott entdecken
            \item Die Eigenschaften Gottes
            \item Ein einziger Gott, drei göttliche Personen
            \item Gott der Schöpfer
        \end{enumerate}      
    
    \end{block}